\textbf{RTP.} Raising bias from .01 to .5 reduced median solve time from .878 to .040 s and median length from 23.46 to 9.39 rad, with 25/25 exact solutions at every setting. Sampling the goal more often helped in this scene. The bias-.5 setting also gave RTP's best median smoothness (.478); clearance instead favored .05 (10.57 cm). Increasing bias did not improve every metric. RTP's full connections from random vertices often produce long detours.

\textbf{RRT.} With bias fixed at .05, increasing range from .2 to 2.5 rad increased the number of exact solutions from 17 to 20 to 24 to 25, but the 2.5 rad setting had a longer median path (8.06 versus 4.12 rad at range 1). At fixed range 1, biases .01/.05/.2 gave 23/24/19 exact trials. Larger bias did not consistently help. Bias .2 had median time .118 s but mean 1.546 s; .05 had .099 s median and .446 s mean. Approximate trials explain much of the tail. Nearest-vertex extension can repeatedly approach obstacles from similar directions, but we did not test this explanation directly.

\textbf{RRT-Connect.} Every setting found exact solutions in all 25 trials. Range 2.5 had the lowest median solve time (.0189 s), though the timing distributions overlap. Bidirectional growth attempts to connect trees from both endpoints, which can reduce exploration needed for these queries. Range .2 gave its shortest median path (4.114 rad); range 5 increased this to 8.827 rad. Median clearance rose from .37 cm at range .2 to 8.83 cm at 2.5, then fell to 6.09 cm at 5. Larger ranges did not always give greater clearance.

\textbf{PRM.} Every neighbor count gave 25 exact solutions out of 25 trials. Median lengths varied only from 8.45 to 8.81 rad; no monotonic improvement followed increasing $k$. The best observed medians were $k=10$ for time (.058 s) and smoothness (.400), and $k=5$ for length and clearance (9.05 cm). PRM starts with a new roadmap in each trial and stops after finding a path.

\textbf{Metric interpretation.} Smoothness sums $[2\theta/(a+b)]^2$ over interior waypoints. Many short edges can score worse than a few long ones: RRT-Connect fell from 557.87 at range .2 to .58 at 5. This measures waypoint geometry, not acceleration. All 478 exact paths passed the prescribed check, but 14 failed OMPL's post-run check with four times the segment count: 4 RRT and 10 RRT-Connect, including 6 at $r=.2$. We kept these paths in the exact-solution summaries because OMPL reported them as exact at the required resolution. That status does not guarantee collision freedom. Sampled minimum clearance can miss collisions. Path-quality summaries use only exact solutions, which may come from different subsets of queries.
